\chapter{Das Hauptfenster}
\label{kap:editor-fenster}

\section{Aufbau}

Das Hauptfenster besteht aus einem festen Mittelteil und mehreren Feldern,
die sich verschieben, abtrennen und ausblenden lassen:

\begin{center}
\begin{tabularx}{\textwidth}{@{}lL@{}}
  \toprule
  Bereich & Inhalt \\
  \midrule
  Schaltplan (Mitte) & Zeichenfläche für die Schaltung oder eine
    Teilschaltung; darüber bei Teilschaltungen eine Leiste mit dem Weg
    zurück zur Hauptschaltung, darunter eine Hinweiszeile mit den Tasten des
    gerade aktiven Werkzeugs. \\
  Text (rechts) & Die Eingabedatei mit Hervorhebung, Zeilennummern und
    Einrückung. \\
  Bauteile & Liste aller Bauteile, Werkzeuge, Teilschaltungen der Datei und
    der Modelle der Bibliothek; mit Suchfeld. \\
  Eigenschaften & Formular für das ausgewählte Bauteil. \\
  Ergebnisse & Analysen, Größen, Wertetabelle und Kurven der letzten
    Simulation. \\
  Protokoll & Ausgabe des Simulators: Fortschritt, Meldungen, Warnungen und
    Fehler, farbig hervorgehoben. \\
  Probleme & Fehler und Warnungen der Prüfung beim Tippen und
    Widersprüche zwischen Plan und Text; Doppelklick springt zur Zeile. \\
  \bottomrule
\end{tabularx}
\end{center}

Die Felder werden an ihrer Titelleiste gezogen. Lässt man eines über einem
anderen los, werden beide zu Reitern; zieht man es aus dem Fenster heraus,
schwebt es frei. \menue{Ansicht} blendet jedes Feld ein und aus. Lage und
Größe aller Felder, die zuletzt geöffneten Dateien und die Einstellungen der
Bibliothek merkt sich der Editor über den Neustart hinweg.

\section{Menüs}

\subsection*{Datei}
\begin{tabularx}{\textwidth}{@{}lP{3.5cm}L@{}}
  \toprule
  Eintrag & Kürzel & Wirkung \\
  \midrule
  Neu & \tasten{Strg}{N} & Neue Datei aus einer Vorlage (RC-Tiefpass mit Plan). \\
  Öffnen \dots & \tasten{Strg}{O} & Eine \datei{.json}-Datei öffnen. \\
  Zuletzt geöffnet & & Die letzten zehn Dateien; \menue{Liste leeren}. \\
  Beispiele & & Die mitgelieferten Beispiele (Anhang~\ref{sec:beispiele}). \\
  Sichern & \tasten{Strg}{S} & Datei sichern. \\
  Sichern unter \dots & \tasten{Umschalt+Strg}{S} & Unter neuem Namen sichern. \\
  Beenden & \tasten{Strg}{Q} & Fragt bei ungesicherten Änderungen nach. \\
  \bottomrule
\end{tabularx}

Eine Datei lässt sich auch beim Start übergeben
(\code{esim-editor datei.json}); unter Windows genügt es, sie auf
\datei{esim-editor.exe} zu ziehen, unter macOS auf das Programmsymbol im
Dock oder über \menue{Öffnen mit}.

\subsection*{Bearbeiten}
\begin{tabularx}{\textwidth}{@{}lP{3.5cm}L@{}}
  \toprule
  Eintrag & Kürzel & Wirkung \\
  \midrule
  Widerrufen / Wiederholen & \tasten{Strg}{Z} / \tasten{Strg}{Y} & Gemeinsam für Plan und Text. \\
  Ausschneiden, Kopieren, Einfügen & \tasten{Strg}{X}, \dots{} & Im Text. \\
  Alles auswählen & \tasten{Strg}{A} & Im Text. \\
  Kommentar ein/aus & \tasten{Strg}{/} & Markierte Zeilen mit \code{//} aus- und einkommentieren. \\
  Auswahl in Teilschaltung umwandeln \dots & \tasten{Strg}{G} & Siehe Abschnitt~\ref{sec:auslagern}. \\
  Modellbibliothek \dots & \tasten{Strg}{L} & Siehe Kapitel~\ref{kap:editor-bibliothek}. \\
  \bottomrule
\end{tabularx}

Unter macOS heißt \menue{Wiederholen} \tasten{Umschalt+cmd}{Z}.

\subsection*{Simulation}
\begin{tabularx}{\textwidth}{@{}lP{3.5cm}L@{}}
  \toprule
  Eintrag & Kürzel & Wirkung \\
  \midrule
  Starten & \tasten{Strg}{R} & Sichert und rechnet die Datei. \\
  Stoppen & \tasten{Strg}{.} & Bricht die laufende Simulation ab. \\
  \bottomrule
\end{tabularx}

Beide Befehle stehen auch in der Werkzeugleiste oben.

\subsection*{Ansicht}
\begin{tabularx}{\textwidth}{@{}lP{3.5cm}L@{}}
  \toprule
  Eintrag & Kürzel & Wirkung \\
  \midrule
  Schaltplan, Text & & Plan oder Text ausblenden (eines bleibt immer sichtbar). \\
  Schaltplan einpassen & \taste{F} im Plan & Ganze Schaltung zeigen. \\
  Bauteile, Eigenschaften, \dots & & Felder ein- und ausblenden. \\
  Kurven: alles zeigen & \tasten{Strg}{0} & Zoom der Kurven zurücksetzen. \\
  \bottomrule
\end{tabularx}

\section{Statuszeile}

Die Statuszeile am unteren Rand zeigt nach jeder Prüfung eine kurze
Zusammenfassung (»3 Bauteile, 3 Unbekannte«) oder die Zahl der Fehler und
Warnungen, während einer Simulation den Fortschritt der laufenden Analyse
in Prozent und danach die Rechenzeit. Auch Rückmeldungen auf Aktionen
(»„NPN\_GEN“ in die Datei übernommen«) erscheinen dort.

\section{Wie der Editor rechnet}

Der Editor startet für jede Simulation das Programm \esim\ als eigenen
Prozess, mit einem temporären Ordner für die Ergebnisse. Das hat zwei
Vorteile: Ein Absturz oder eine Endlosschleife des Simulators trifft nicht
den Editor, und \menue{Stoppen} beendet einfach den Prozess. \esim\ liegt im
Programmpaket neben dem Editor.

Die Prüfung beim Tippen läuft dagegen im Editor selbst, mit demselben
Parser wie der Simulator, in einem Hintergrund-Thread und kurz verzögert,
damit sie das Tippen nicht bremst. Was sie meldet, meldet auch der
Simulator.

\section{Start mit Optionen}

\begin{shell}
esim-editor [datei.json] [--subcircuit <name>] [--run] [--snapshot <bild.png>]
\end{shell}

\begin{tabularx}{\textwidth}{@{}lL@{}}
  \toprule
  Option & Wirkung \\
  \midrule
  \code{datei.json} & Öffnet die Datei. \\
  \code{--subcircuit name} & Zeigt den Plan der Teilschaltung \code{name}. \\
  \code{--run} & Startet gleich die Simulation. \\
  \code{--snapshot bild.png} & Sichert das Fenster (nach der Simulation) als Bild und beendet sich. Gedacht für Dokumentation und Tests; alle Bildschirmfotos dieses Handbuchs sind so entstanden. \\
  \bottomrule
\end{tabularx}

Unter macOS liegt das Programm im Paket:

\begin{shell}
esim-editor.app/Contents/MacOS/esim-editor
\end{shell}

\section{Hell und dunkel}

Der Editor folgt der Darstellung des Betriebssystems. Plan, Kurven und
Symbole haben für den hellen und den dunklen Modus eigene, aufeinander
abgestimmte Farben. Die acht Kurvenfarben sind so gewählt, dass sie auch bei
Farbfehlsichtigkeit unterscheidbar bleiben; Kurvenscharen aus
Parameterstudien erscheinen als Stufen einer Blauskala.
